% Narrative literature review template (LaTeX)
% This is a lightweight starting point for surveys/related-work sections.

\documentclass[11pt]{article}

\usepackage[T1]{fontenc}
\usepackage[utf8]{inputenc}
\usepackage{lmodern}
\usepackage{microtype}
\usepackage{amsmath,amssymb}
\usepackage{graphicx}
\usepackage{booktabs}
\usepackage{hyperref}
\usepackage[numbers]{natbib}

\title{Title of Literature Review}
\author{Author Name}
\date{\today}

\begin{document}
\maketitle

\begin{abstract}
State the scope of the review, the selection methodology (briefly), and the main themes and gaps (150--250 words).
\end{abstract}

\section{Introduction}
Define the problem area and why it matters. Clarify the scope and what is out-of-scope.

\section{Method (Source Selection)}
Describe where you searched (databases), the time window, and the inclusion/exclusion criteria. This is required for defensible reviews.

\section{Thematic Synthesis}
Organize work by themes rather than by paper-by-paper summaries.

\subsection{Theme A}
Summarize what is known and where results agree or disagree \citep{smith2023,jones2022}.

\subsection{Theme B}
Compare approaches under consistent dimensions (assumptions, data, metrics, settings).

\section{Open Problems and Gaps}
List gaps backed by evidence (e.g., missing ablations, weak external validity, lack of reproducibility artifacts).

\section{Conclusion}
Summarize themes and propose future directions.

\bibliographystyle{plainnat}
\bibliography{references}

\end{document}

